\FloatBarrier
\section{Additional Work}
\label{sec:hybrid-additional-work}

Small-World Allocation (SWA) examines whether functional organization among attention heads can help characterize and preserve language-model capability. Inspired by functional brain-network analysis, the study builds graphs over attention heads using singular vector canonical correlation analysis of activations on fluid-reasoning questions. Across six open-source models, higher small-world indices were associated with lower answer-token negative log-likelihood on DRE-Bench. To translate this graph-level observation into a practical method, SWA derives core scores for within-community connectivity and bridge scores for connectivity across communities. These scores guide hierarchical sparsity allocation across Transformer blocks and individual attention heads, preserving components with stronger structural roles. Experiments combined this allocation with SparseGPT and Wanda and evaluated language-model perplexity, zero-shot accuracy, and held-out reasoning tasks. The results showed improved model performance under pruning, while functional-graph measurements showed better preservation of small-world organization. Ablations indicated that layer-level allocation supplies the main signal and head-level allocation further improves performance, especially at higher sparsity. As an additional study alongside \hybrid{}, SWA extends the examination of performance-relevant relational structure from brain regions to components of artificial neural networks. \hybrid{} learns higher-order brain relationships for cognitive prediction, while SWA uses functional network organization to guide the retention of computational structure.
